\chapter[Polar Frame and Lorentz Read]{The Polar Frame and the Local Lorentz Read}
\label{chap:pair-production-geometry-generating-gauge}

\section{The Local Frame}

The previous chapter left the grammar with one concrete reader for the
positive invariant. The invariant was no longer a floating name. It had an
anchored operator, a scale, a kernel of labels, and a rule for deciding when a
sensor class could carry a reading. The present chapter asks what that
reader is forced to do when the first charged difference appears. The answer
cannot be imported as a particle story. It must be read from the grammar that
has already been earned.

The pair problem begins at the point where stationarity is too poor to
finish the record. The first variation has vanished. The null drift has been
pinned. The second variation is the surviving activity. But a quadratic
activity can be read in more than one face. It can be read as magnitude. It
can be read as strain. It can also be read as an oriented charge when the
baseline against which the activity is compared has been fixed. Pair
production is the grammar's name for the moment when that oriented reading
cannot be avoided.

This is why the chapter does not begin with two preexisting objects. A pair
is not placed into the ledger and then assigned signs. The signs are forced
by the way the ledger reads a coupled disturbance. If the coupled residue
telescopes away, no charged distinction remains. If the residue refuses to
cancel under the permitted action, the grammar must discriminate it. The
discrimination is signed because it is read relative to a baseline. The
baseline is not decorative. It decides which orientation is native and which
orientation appears only after a tilt.

The chapter moves through seven obligations. First, the electron emerges as
the flat-baseline reading of the discriminating action, and the charge is
\(-1\). Second, the positron appears only when the same pair is read against
a tilted baseline, where the sign is forced to \(+1\). Third, the chapter
isolates the two identity grants on which the physics rests: the local
quotient that lets selection stand after collapse, and the motion grant that
lets stationarity answer to the world. Fourth, boundary geometry generates
the gauge by producing the action that discriminates the pair residue. Fifth,
the geometric and gauge readings split; a finite carrier may read one side or
the other, but it may not form an interior direct sum that hides a mixed
scalar. Sixth, the electromagnetic reading becomes the Cauchy limit of the
geometric reading as the mesh residue is forced to zero. Seventh, finite
Cohen forcing names the disciplined ledger of bounded spline conditions by
which the continuum reading is compelled without being assumed.

Examples will appear, but only as lanterns. A coincidence detector may
illustrate the cancellation of an electron and a positron. A loop holonomy
may illustrate why an open path carries no charge while a closed path carries
a sign. A compact disc reader may illustrate a threshold gauge, and the
Cantor--G\"odel--Cohen warning may illustrate why completion cannot be
accepted without recoverability. None of these examples supplies the spine.
The spine is the grammatical obligation: what the finite reader permits,
what it forbids, what it identifies, what it quotients, what it carries, and
what it is forced to read at the boundary.

The title therefore names one process, not two. Pair production and geometry
generating gauge are the same turn seen from opposite sides. From the charge
side, the second variation discriminates a pair and forces the electron
reading. From the boundary side, geometry is not a container but the
stabilized grammar of separations; when that boundary grammar is asked to
carry the discriminating residue, it generates the gauge that reads the pair.
The chapter's task is to keep those two faces together without collapsing
them into an illicit direct sum.

\section{The Signed Norm}

The electron enters only after the grammar has exhausted cheaper readings.
An unanchored drift cannot be charged, because the boundary has not yet
forbidden it. A first residual cannot be charged at stationarity, because the
stationary reader has forced it to vanish. A merely positive activity cannot
yet be an electron, because magnitude alone has no orientation. The electron
appears when the second variation is read by an action that discriminates the
coupled residue rather than letting it telescope into silence.

The discriminating action is the action that refuses the wrong cancellation.
A telescoping action reads adjacent debts so that the left contribution is
answered by the right and the boundary closes without a signed remainder.
Such an action is necessary in the earlier chapters because it explains how
stationarity can be certified. But the pair problem is different. A pair kick
couples two readings that cannot both be absorbed as independent left and
right debts. The grammar must ask whether the coupled difference is merely
the sum of its separated parts or whether a residue remains after those
parts have been accounted for.

The electron is the remaining reading. The grammar compares the coupled
cubic difference with the left and right readings taken separately. If the
coupled reading were no more than those separated readings, the pair would
not carry a new charge. But the mixed residue is not permitted to vanish by
renaming. It survives the subtraction of the separated parts. Because the
first variation has already been driven to zero, this surviving residue is
read at second order. Because the baseline is flat, the orientation of that
second-order residue is fixed. The grammar reads the charge as \(-1\).

This sign is not an ornament attached to a particle after the fact. It is the
orientation of the discriminated residue. The flat path supplies the native
baseline: the comparison in which no tilt has been introduced, no loop has
changed the orientation, and no external convention has reversed the sign.
Relative to that baseline, the coupled residue points against the positive
orientation. The grammar therefore reads electron, not because a name has
been announced, but because the only admissible charged reading of the flat
pair is \(-1\).

The word charge should be handled with the same discipline as the word
energy in the previous chapter. Energy became readable only after the
anchored operator could detect zero activity. Charge becomes readable only
after the baseline can discriminate orientation. A magnitude without a
baseline has size but no sign. A sign without the second variation has
orientation but no activity to carry it. The electron combines the two: a
second-variation activity whose discriminating action reads negative
orientation over the flat path.

This also explains why the electron emerges from the invariant rather than
beside it. The grammar has not introduced a second measure. It has taken the
single positive invariant and asked how it reads when the pair residue cannot
be decomposed into separated debts. The invariant supplies the activity; the
baseline supplies orientation; the discriminating action supplies the refusal
of false cancellation. The result is a charged reading of the same invariant.

The flatness of the path matters because it forbids the positron reading at
this stage. Over the native baseline, a \(+1\) charge would require either a
tilt in the comparison or a reversal of the sign convention. Neither is
available inside the flat reading. The grammar therefore does not merely
permit the electron reading. It forbids the opposite reading until the
baseline changes. The electron is forced because every other charged
interpretation would spend a distinction the flat grammar has not granted.

A finite coincidence example can illuminate the sign without supplying the
argument. When a detector later reads two back-to-back photons inside a
finite energy window, it treats a prior electron-positron encounter as a
cancellation of signed charges. The cancellation is visible only because the
opposed signs have already been distinguished. In the present section the
grammar is earlier than that detector story. It is establishing the native
negative sign that such a later cancellation can use.

Thus the electron is not an assumed inhabitant of the ledger. It is the
grammar's first charged reading of the single invariant. The second
variation carries activity. The pair kick exposes a mixed residue. The
discriminating action refuses to quotient that residue away. The flat
baseline fixes the orientation. Those obligations close on the charge
\(-1\), and the name electron is the reading of that closure.

\section{Radial, Polar, and Azimuthal Legs}

Once the electron has been read over the flat baseline, the opposite sign
cannot be obtained by wish. The grammar has already forbidden it there. A
positron reading requires a different comparison, not a second arbitrary
label. The same pair must be carried against a tilted baseline. Only then
does the charged residue face the reader with the opposite orientation, and
only then is the sign \(+1\) admissible.

Baseline dependence is not a weakness in the charge reading. It is what
makes the sign meaningful. A sign is not a magnitude; it is a relation to an
oriented reference. The flat path supplies one reference. The tilted path
supplies another. If the grammar erased the difference between them, it would
erase the very orientation it is trying to read. A positron is therefore not
the electron plus a new substance. It is the same pair grammar seen through a
comparison whose baseline has been tilted.

The tilt is a constrained act. It does not let the reader reverse any sign at
will. A permissible tilt must preserve the underlying pair, the second
variation, the boundary commitments, and the distinction between separated
parts and coupled residue. What changes is the orientation of the baseline
against which the residue is read. Under that changed comparison, the same
discriminated activity that read \(-1\) over the flat path reads \(+1\). The
grammar identifies the sign flip as baseline-relative, not as a second
independent event.

This is why the native baseline carries an asymmetry. Over the flat path the
electron is present and the positron is not. The positron appears only when
the reading has been tilted enough to make the opposite orientation
admissible. The asymmetry is not imposed by a story about preference. It is
forced by the grammar of what the baseline can carry. Native flat reading
permits the negative charge and forbids the positive one. Tilted reading
permits the positive charge by changing the orientation of comparison.

Loop language helps, provided it remains an illustration. An open path can
carry local comparisons and still return no charge, because it has not
closed the orientation into a recoverable holonomy. A closed loop may return
with a signed residue, even when local segments looked balanced. The
important point is not that every positron must be pictured as a loop. The
point is grammatical: sign appears when transport around the relevant
comparison closes with an orientation that the reader cannot quotient away.
The tilted path supplies that closure for the \(+1\) reading.

The Aharonov--Bohm example shows the same form at a later, physical register.
Two branches can travel through regions where the local field reading is
silent, then recombine with a phase difference because the connection carried
memory around the loop. The present chapter uses the example only to sharpen
the grammar: local flatness does not guarantee global triviality, and the
sign of a transported distinction can depend on the path by which the
baseline is closed. The positron is read when the path carries the opposite
orientation.

Annihilation gives a second illustration. A coincidence detector accepts two
opposed photon directions inside a finite window as evidence that opposite
charges have cancelled. That finite coincidence logic presupposes two signs:
the electron's \(-1\) and the positron's \(+1\). The detector does not create
the signs. It reads the cancellation that those signs permit. In the grammar
of this chapter, the positron sign has already been forced by the tilted
baseline before any later coincidence story can use it.

The crucial discipline is that the positron reading is neither arbitrary nor
available everywhere. If a proposed reading announces \(+1\) over the flat
path, it violates the baseline. If it announces \(+1\) after a permissible
tilt, it is not inventing a new charge; it is reading the same
second-variation activity under the opposite orientation. The grammar
therefore preserves both facts: the electron is native to the flat baseline,
and the positron is the tilted-baseline charge.

The pair production grammar now has a real split in sign. The signs cancel
when read together, but they are not indistinguishable before cancellation.
The flat baseline reads the electron. The tilted baseline reads the
positron. A reader that cannot carry the tilt cannot read the positron as a
physical record; it must treat the positive sign as unavailable or
metaphysical relative to that class. A reader that can carry the tilt reads
the opposite orientation without changing the underlying invariant.

Thus the positron is not a second foundation. It is the charged reading
forced by a change in baseline. The same grammar that forbids \(+1\) over
the flat path permits \(+1\) over the tilted path. Baseline dependence is the
condition of the reading, and the sign flip is the grammar's first clear
lesson that orientation belongs to the relation between invariant and
reader.

\section{The Orientation Convention}

The chapter can now state how little it has used. The electron and positron
readings do not require a new stock of primitive objects. They rest on the
two grants already forced by the earlier grammar. The first grant is the
local identity that survived the collapse of truth: the needle, or selection
quotient, by which a reading may stand as the same for a specified purpose.
The second grant is the motion grant: the permission to let the world answer
the stationarity question by fixing which residuals vanish under admissible
variation.

Calling these grants axioms is safe only if the word remains narrow. An
axiom here is not a license to import a hidden universe. It is a bounded
permission without which the grammar cannot continue. The first chapter
showed why unbounded identity was forbidden. The second chapter showed that
finite variation needs a rule by which stationarity can be read. The third
chapter grounded the invariant in a concrete operator. The present chapter
uses exactly those permissions and no third one.

The first grant identifies what would otherwise remain only selected. A
quotient does not say that all differences have vanished. It says that, for
this reading, certain differences no longer count as distinct. The charge
reading needs this grant because a baseline is a selection: it identifies the
reference against which orientation is read. Without the local quotient, the
flat path and the tilted path would be mere names with no stable relation to
the pair residue. With the quotient, each baseline can carry its own
identity while still being distinguished from the other.

The second grant is different. It is not a quotient of presentations but a
world-facing constraint on motion. The grammar may define permitted
variations, actions, boundaries, and residuals, but it still needs the
reading in which stationarity is accepted as the law of the situation. This
grant says that a vanishing residual is not merely a formal convenience; it
is the condition by which the measured path answers. It is the place where
the grammar asks the world for the rule of motion and then carries the
answer as constraint.

The two grants are enough because the charged readings use no additional
privilege. The electron needs the selection of a flat baseline and the
stationary second-variation grammar. The positron needs the selection of a
tilted baseline and the same stationary grammar. Geometry generating gauge
will need boundary identity and motion constraint, not a new substance.
The split will need a carrier discipline that forbids illegal mixing, not a
third metaphysical key. The Cauchy limit and finite forcing will need
extension, compatibility, and residual bounds, all of which belong to the
finite grammar already earned.

This austerity matters because signs are tempting. Once a negative and a
positive reading appear, the reader may want to populate a whole world of
opposed entities by declaration. The grammar forbids that. A sign must be
carried by the invariant, oriented by a baseline, and admitted by a sensor
class. A name that lacks those supports is not yet a reading. The two grants
keep the book from replacing measurement with mythology.

The first grant also prevents the opposite mistake, in which no identity is
allowed and therefore no charge can be read. If the flat baseline could not
be selected as the same baseline across the comparison, the \(-1\) reading
would dissolve. If the tilted baseline could not be selected as its own
comparison, the \(+1\) reading would never become stable. The quotient is
therefore not a philosophical luxury. It is the grammar's way of permitting
orientation to have a reference.

The second grant prevents a merely conventional sign system. If the world
did not answer the stationarity question, a sign could be assigned but not
measured. The charged residue would be a mark in a ledger without a rule
that says which variations are admissible and which residuals must vanish.
The law-of-motion grant bounds the possible readings. It forbids a sign from
being detached from the action that discriminates it.

Together the grants explain why the chapter can be bold without becoming
loose. It reads electron and positron, but it reads them as forced signs of a
single invariant under specified baselines. It lets geometry generate gauge,
but only through boundary grammar. It reaches a continuum reading, but only
through finite conditions. The two grants are the only doors through which
identity enters the chapter. Everything else must be carried, bounded,
quotiented, or forced.

\section{Geometry Generates Gauge}

Geometry enters this chapter as boundary grammar, not as a container waiting
behind the record. A geometry is a stabilized way of saying which separations
matter, which boundaries have been committed, which interiors remain
unread, and which transports preserve distinguishability. When such a
boundary grammar is asked to carry the charged pair residue, it cannot remain
merely spatial. It must generate a gauge: a rule for reading how the
distinction is transported, compared, and discriminated.

The generating step begins with a boundary depth. An interior reading may
carry no charge when read alone. A boundary reading may carry a unit
obligation because it is where the interior must answer to an exterior
comparison. The grammar distinguishes these two faces. It forbids the
interior from inventing charge without boundary support, and it forbids the
boundary from acting as a decorative shell. The boundary is the place where
unspent orientation becomes readable.

A compact-disc threshold gives a modest illustration. The read head does not
begin with a metric as a metaphysical object. It admits a new mark only when
two surface configurations are separated enough to force a new event in the
record. The gauge of separation is inferred from admissible distinction. In
the chapter's grammar, spatial boundary works similarly. A boundary is not a
line drawn around a pre-given space. It is the rule that says when a
separation has become expensive enough to enter the ledger.

Once the boundary carries this rule, it can be compiled into an action. The
word compiled should be read grammatically: boundary commitments are gathered
into a finite reader that compares interior and exterior charges, preserves
the anchor, and tests the coupled residue. The resulting action is not added
beside geometry. It is geometry's boundary obligation expressed as a gauge
test. That test is precisely what the electron section called the
discriminating action.

This is the sense in which geometry generates gauge. The grammar starts with
boundary separations and asks what must be carried so that a pair residue can
be read without contradiction. It must carry a baseline. It must carry
orientation. It must carry the distinction between separated parts and the
coupled residue. It must carry the boundary debt that the interior alone
cannot settle. Those carriers are gauge carriers. The gauge is the transport
grammar forced by the boundary.

The Aharonov--Bohm illustration again shows the shape. A region may have no
local field reading along two branches, yet the closed comparison can return
a phase difference. The physically relevant distinction is then not the
local force reading alone but the connection that carries memory around the
loop. In this chapter the point is more basic: a boundary grammar that
preserves distinguishability around a comparison must become a gauge grammar
when the transported sign matters.

This conversion also protects the electron reading from arbitrariness. If
the discriminating action were merely chosen, the \(-1\) reading would look
like a convention. But when boundary geometry generates the action, the
charge reading is forced by the boundary's own obligation to reconcile
interior and exterior. The gauge reads what the geometry has made
unavoidable. It discriminates because the boundary has produced a residue
that cannot be quotiented away.

The white-hole image can be used carefully here. A source boundary exports a
reading that the interior cannot absorb as its own private fact. The export
does not mean a mythic object has been placed at the boundary. It means the
boundary grammar carries an asymmetry: the interior may be charge-silent,
while the boundary carries the unit condition needed for discrimination.
Gauge is the rule that transports that boundary condition into a readable
action.

Thus geometry generating gauge is not a slogan of unification. It is a
discipline of origin. The grammar permits geometry as stabilized boundary
bookkeeping. It forbids geometry from claiming charged transport until a
residue demands it. It then identifies the necessary transport rule as gauge.
When the gauge acts on the pair residue, the electron reading is not an
extra particle inserted into geometry. It is the boundary-generated action's
negative charge.

The section also prepares the split. If gauge is generated by geometry, it
is tempting to add the geometric and gauge readings together as if they were
two independent components of one finite vector. The grammar will forbid
that. Generation is not direct summation. The gauge is born from the boundary
obligation of geometry, but once the reading is made, the finite carrier must
respect the difference between the geometric face and the gauge face.

\section{The Split}

The geometry-to-gauge conversion creates a new danger. Because the gauge is
generated by boundary geometry, the reader may try to carry the geometric
sector and the gauge sector as two independent summands. The grammar forbids
that finite direct sum. A finite certificate may read the geometric face or
the gauge face, but it may not form an interior mixed scalar that pretends
both faces are simultaneously available as independent components.

The split is therefore a carrier discipline. It says how a reading may pass
through the Hilbert door without losing the distinction between the
geometric source and the gauge sign. The geometric carrier reads boundary
separation, curvature-like residue, and the mesh on which the ledger is
refined. The gauge carrier reads transported orientation, charge, and the
discriminating action. The same invariant underlies both, but the finite
reader is not allowed to add them as if it possessed a larger interior space
where their mixed product were already meaningful.

The reason is the mixed term. If geometry and gauge were freely summed, an
interior scalar cross-reading would be available. That cross-reading would
say that the geometric residue and the gauge residue can be multiplied
inside the same finite certificate without remainder. But the chapter has
not earned such a certificate. The boundary generated the gauge precisely
because the interior could not settle the charged residue on its own. To
reinsert the gauge as an interior summand would erase the boundary
obligation that generated it.

Orthogonality is the disciplined alternative. It does not mean the two
sectors are unrelated. It means the finite grammar reads them in separated
channels whose interior scalar mixing is forbidden. The geometric channel
can carry the boundary-to-interior refinement. The gauge channel can carry
the sign transport. What cannot be carried is a hidden interior scalar that
lets one channel cancel debt in the other. The split protects the charged
reading from being dissolved into geometry and protects geometry from being
reduced to a gauge label.

This explains why boundary radiation is waiting in the next chapter. If the
mixed term cannot live in the interior, then any reconciliation between the
geometric and gauge faces must appear at the boundary. Radiation will be the
grammar's name for that boundary phenomenon. The present chapter does not
need to derive it yet. It only has to force the condition that makes it
unavoidable: no interior mixed certificate.

The split also clarifies matter and antimatter readings. Matter and
antimatter are not two substances thrown into a common box. They are signs
read through the gauge under specified baselines. The native flat baseline
reads the electron side. A tilted baseline reads the positron side. If the
geometric and gauge sectors were collapsed into a direct sum, the baseline
dependence would blur. The grammar would no longer know whether a positive
sign came from a tilt, a boundary, or an illicit interior mixture. The split
keeps the sign honest.

The Yang--Mills illustration helps at the level of labels. Different
refinement classes may require local repair rules that close under their own
composition and commute only where their labels are independent. The smooth
shadow may look like a product of symmetries, but the ledger discipline is
more severe: independent label classes must not corrupt one another's
bookkeeping. The present split is even sharper. Geometry and gauge are not
two free label classes to be added in the interior. They are two faces of one
boundary-forced reading, and their mixed scalar is unavailable.

The split is not a failure of unification. It is the form unification takes
under finite measurement. A loose unification would announce one large space
and place everything inside it. The grammar refuses. It identifies the
single invariant underneath both readings, but it also preserves the carrier
discipline that makes each reading meaningful. The unity is in the invariant
and the boundary process. The split is in the finite certificates by which
the readings are carried.

Thus the geometry/gauge relation is neither separation without contact nor
fusion without discipline. Geometry generates the gauge at the boundary.
Gauge reads the charged residue. The finite Hilbert pairing door then forces
a split carrier: geometric or gauge, not an illicit sum. This is why the
chapter can say that the sectors do not combine as a direct sum while still
saying that one grammar has forced both faces.

\section{The Zero-Mesh Reading}

The split leaves a question of scale. If geometry generates gauge and the
finite carrier keeps the two faces apart, how can the electromagnetic
reading be related to the geometric reading without forming the forbidden
interior sum? The answer is Cauchy convergence. The electromagnetic reading
is the zero-mesh limit of the geometric reading: not a completed continuum
assumed in advance, but the limiting gauge face forced as geometric
residuals are refined toward zero.

Begin with the geometric ledger. It carries a mesh, boundary separations,
and residuals that measure how much geometric refinement has not yet been
absorbed. A coarse mesh reads a coarse boundary grammar. A finer mesh can
split cells, add anchors, and test narrower supports. At each stage the
geometric reading is finite. It has not inspected every point. It has only
committed the distinctions its mesh can carry.

The gauge reading is what remains stable under the refinement of that mesh.
As the geometric residual is bounded by a visible \(\epsilon\), and as
\(\epsilon\) is forced downward, the transported orientation becomes less
dependent on the coarse presentation. The charge reading does not wait for a
completed continuum to arrive from outside. It is the value that every
admissible refinement is being squeezed toward. In this sense the
electromagnetic reading is the Cauchy limit of the geometric reading.

The shape is the same as the earlier finite approach to real value:
\[
  0 \leq R_n \leq \epsilon_n,
  \qquad \epsilon_n \to 0 .
\]
Here \(R_n\) is the unread geometric residue of the \(n\)-th mesh, and the
gauge face is read when the residual sequence becomes Cauchy under the
allowed refinements. The lower bound forbids hidden negative debt from
cancelling a positive residue. The upper bound records the mesh error still
being carried. The vanishing of the upper bound forces the limiting reading.

This is what the chapter means by saying that QED is the Cauchy limit of GR.
The phrase should not be inflated into a claim that every smooth physical
theory has been derived from every other. It names a specific grammatical
relation. The geometric reading carries mesh, boundary, and curvature-like
residue. The electromagnetic reading carries the gauge sign that remains as
the geometric mesh is forced to zero. The two are related by a Cauchy
process, not by an interior direct sum.

The Precision example gives the right caution. A spline completion may
assign determinate values between anchors, but those intermediate values are
not new measurements. They are consequences of the admissible completion.
Likewise, the zero-mesh gauge reading is not an observed point inside a
completed continuum. It is the compelled reading of a refinement process
whose residuals have been bounded and forced downward. Precision is a
consequence of coherence; accuracy remains answerable to future records.

This section also explains why the split was necessary. If the geometric and
gauge sectors could be added in the interior, the limit would be too cheap.
One would simply place the electromagnetic term beside the geometric term
and call the result unified. The grammar instead requires a refinement
history. Geometry carries the mesh. Gauge carries the transported sign. The
limit identifies the stable sign as mesh residue vanishes. No interior mixed
scalar is introduced.

The zero-mesh reading is therefore an exterior discipline on a finite
process. Each mesh stage is finite. Each residual bound is finite. Each
extension must preserve previous commitments or explicitly refine them. The
Cauchy property says that after sufficient refinement, later stages cannot
be distinguished by the permitted residual tests beyond the chosen bound.
The electromagnetic reading is the quotient of those late-stage readings
that no admissible test can tell apart.

This is also why baseline dependence survives the limit. Refinement does
not erase orientation. It squeezes geometric residue while preserving the
baseline relation that makes sign readable. A flat-baseline refinement tower
continues to read the electron sign. A tilted-baseline tower continues to
read the positron sign. The limit does not abolish the grammar of
orientation; it makes the transported orientation independent of mesh
presentation.

The chapter has now moved from charged signs to boundary-generated gauge,
from gauge to split carriers, and from split carriers to a limiting
relation. One obligation remains. The Cauchy limit requires a disciplined
family of finite conditions. The grammar must say how those conditions are
ordered, extended, made compatible, and forced to decide the continuum
reading. That is the role of finite Cohen forcing.

\section{Finite Cohen Forcing}

The continuum reading at the end of the previous section cannot be accepted
as a gift. A Cauchy process needs finite conditions that make its approach
honest. Finite Cohen forcing is the chapter's name for that discipline. It
does not import an unbounded set-theoretic machine. It names a finite,
\(\epsilon\)-bounded grammar of spline atoms, supports, extensions,
compatibility tests, and dense decisions. The continuum reading is forced
because every admissible finite condition can be extended until the relevant
question is decided within the bound.

A spline atom is the smallest local completion unit the chapter needs. It
has an anchor, a bounded region of support, and a value or activity reading
that can be compared with neighboring atoms. An atom by itself is not a
continuum. It is a finite obligation: here is a local piece of the ledger,
here is what it carries, here is the bound under which it may be extended.
A spline condition is a finite collection of such atoms arranged so that
their overlaps are compatible.

Finite support is essential. If a condition were allowed to carry infinitely
many active obligations at once, it would no longer be a measurement record.
It would pretend to have spent all future refinements in advance. Finite
support says that only finitely many atoms are currently active. Other atoms
may be added later by extension. The grammar therefore preserves the
halting horizon even while it points toward completion.

Extension is the order of forcing. A stronger condition does not abandon the
weaker one. It preserves the commitments already made and adds refinements,
smaller bounds, or more local decisions. If a proposed extension contradicts
an earlier atom on their shared support, it is not an extension at all. It is
a different ledger. Compatibility is the grammar's refusal to let future
precision erase past record without paying for the revision.

Dense decision is the forcing heart. For every question the grammar is
authorized to ask, and for every current condition that has not answered it,
there must be an admissible extension that decides it up to the current
\(\epsilon\). The word dense means that undecided conditions are not allowed
to sit in isolated pockets forever. Wherever the finite ledger stands, the
question can be pressed further by an extension. The decision may be
positive, negative, or bounded as unreadable, but it must be made in the
language of the finite condition.

This is how the continuum is forced rather than assumed. The ledger does not
survey all possible points. It builds a chain of finite conditions such that
every required local question is eventually decided within a shrinking
bound. Two completed readings are identified when no finite dense decision
can separate them. A proposed continuum distinction that cannot be recovered
by any such finite extension is forbidden as unearned structure.

The Cantor--G\"odel--Cohen example serves as a warning. Different
completions may agree on the finite records a grammar has chosen to carry.
Consistency alone is therefore too weak. A completion is admissible here
only when its distinctions are recoverable from finite refinement. Finite
Cohen forcing gives the grammar a way to say which completions are forced by
the ledger and which are optional stories placed above it.

The compact-disc threshold example gives the same warning in concrete form.
A coarser reader may record fewer events than a finer reader. When the
records are merged, the finer distinctions force refinements on the coarser
history where compatibility requires them. But distinctions below every
available threshold do not enter the record merely because a continuum could
name them. They must be forced by an admissible extension of the ledger.

For Chapter 04, the forced continuum reading is the stable home of the
pair-production grammar. The electron sign, the tilted positron sign, the
boundary-generated gauge, the split, and the zero-mesh relation all require
finite decisions to survive refinement. A condition may decide a local sign.
Another may decide a boundary compatibility. Another may shrink a geometric
residual. The forcing process keeps extending until every required decision
is dense in the refinement order.

This does not erase finitude. It honors it. The continuum reading is the
quotient of refinement histories that finite dense decisions can no longer
distinguish. The grammar carries atoms, bounds, supports, and extensions. It
forbids unrecorded intermediate structure. It identifies compatible
late-stage readings. It forces the continuum as the only reading left after
all finite questions have been pressed under shrinking bounds.

The chapter has therefore reached a full circle. It began with the concrete
second variation and read the electron. It tilted the baseline and read the
positron. It isolated the two grants that make those readings lawful. It let
boundary geometry generate the gauge. It split geometry and gauge to forbid
interior mixing. It read the electromagnetic face as the Cauchy limit of the
geometric face. Finally, it forced the continuum reading by finite
conditions. The next chapter inherits the consequence: when the interior
mixed term is forbidden and the continuum reading is forced, radiation must
be read at the boundary.

\section{What the Local Frame Carries Forward}

The chapter has kept the single invariant and made it charged. The electron
was read as the \(-1\) value of the discriminating action over the flat
baseline. The positron was read as the \(+1\) value of the same pair grammar
only after the baseline tilted. The two signs did not arrive as free
opposites. They were forced by the relation between second variation,
baseline, and discriminated residue.

The chapter also kept identity narrow. It used the local quotient that lets
a selected baseline stand as the same baseline for a specified reading, and
it used the motion grant that lets stationarity answer to the world. It did
not add a third source of identity. Under those two grants, boundary geometry
generated the gauge by producing the action that reads the charged pair.
The gauge did not replace geometry, and geometry did not swallow the gauge.
Their finite carrier split.

That split is what the next chapter receives. Since the geometric and gauge
faces do not form an interior direct sum, their reconciliation cannot be an
interior scalar mixed term. The grammar must look to the boundary. The
zero-mesh Cauchy reading and the finite Cohen forcing ledger have supplied
the continuum discipline needed for that next move. What remains is to read
how the boundary carries the residue that the interior is forbidden to mix.

Chapter 05 begins there. The present chapter has forced the charged pair,
the baseline dependence, the geometry-to-gauge conversion, the split, and
the finite forcing of the continuum reading. Boundary radiation is not an
extra topic after those results. It is the next grammatical obligation.
